\chapter{相关工作}
\label{chap:related}

本章回顾与本工作相关的四个方向：多模态融合方法、老年人跌倒检测、环境声音分类，以及本工作与现有研究的差异定位。

\section{多模态融合方法}

多模态融合按层次可分为早期融合（特征级）、晚期融合（决策级）和中间融合（模型级）\cite{baltrusaitis2019multimodal}。简单的特征拼接（concatenation）属于早期融合，无法建模模态间的交互关系。基于注意力的融合机制能够学习模态间的相关性，是中间融合的代表。

\subsection{Transformer与跨模态注意力}

Vaswani等\cite{vaswani2017attention}提出的Transformer架构以自注意力为核心，能够捕获序列内长距离依赖。该机制自然延伸到跨模态场景：不同模态的特征作为Query/Key/Value，通过注意力计算实现模态间信息交互。

ViLBERT\cite{lu2019vilbert}通过联合遮蔽语言建模和视觉区域预测任务预训练视觉-语言模型，验证了跨模态注意力在多模态对齐中的有效性。ALBEF\cite{li2021align}进一步提出``先对齐再融合''的策略，使用动量蒸馏提升视觉-语言表示学习。这些工作证明了跨模态注意力机制能够有效建模异构模态间的交互关系。

本工作将Transformer的跨模态注意力思想应用于老年人监护的四模态融合场景，使每个模态能够attend到所有其他模态，实现联合预测。

\subsection{固定权重与学习权重的融合}

传统的多模态融合常采用固定权重加权（如视频0.4/音频0.25/生理0.25/用药0.1），权重由专家经验设定。这种方式的局限在于：（1）权重无法随输入自适应调整；（2）无法建模模态间的细粒度交互。

本工作采用可学习的模态权重（经softmax归一化）与跨模态注意力相结合，使模型既能学习全局模态重要性，又能通过注意力机制实现细粒度的跨模态交互。

\section{老年人跌倒检测}

跌倒检测方法可分为三类\cite{rachakonda2020cstick}：

\begin{enumerate}
    \item \textbf{基于穿戴传感器}：利用加速度计/陀螺仪检测跌倒时的加速度峰值。精度较高但需老人主动佩戴，依从性差。
    
    \item \textbf{基于视频的姿态分析}\cite{auvinet2010fall}：通过摄像头捕获人体姿态，分析重心下降和躯干角度变化判断跌倒。非接触式但受光照、遮挡影响。
    
    \item \textbf{基于音频的撞击声识别}\cite{popescu2009acoustic}：通过麦克风捕获跌倒撞击声，利用声学特征分类。成本低但易受环境噪声干扰。
\end{enumerate}

单一方法均有局限，多模态融合是提升检测率与降低误报率的有效途径。本工作融合了视频姿态分析（MediaPipe骨架）和音频声学特征，辅以生理数据和用药信息，实现更全面的风险评估。

\subsection{MediaPipe姿态估计}

MediaPipe\cite{lugaresi2019mediapipe}是Google开源的跨平台机器学习推理框架，其Pose解决方案可实时检测33个人体关键点。相比OpenPose等模型，MediaPipe轻量且支持移动端部署。在MediaPipe 0.10.35版本中，旧的\texttt{mp.solutions.pose} API已被移除，迁移至Tasks API的\texttt{PoseLandmarker}，提供了更稳定的关键点检测和更灵活的运行模式配置。

本工作采用MediaPipe Tasks API提取骨架特征，利用33个关键点的时序信息（中心轨迹、躯干角度、运动速度）构建视频特征向量。

\section{环境声音分类}

ESC-50\cite{piczak2015esc}是环境声音分类的标准数据集，包含2000条5秒音频，涵盖50个类别（动物声、自然声、人声、室内声、室外声）。基于librosa\cite{mcfee2015librosa}提取的MFCC、Mel频谱等声学特征是音频分类的经典方法。

Audio Spectrogram Transformer (AST)\cite{gong2021ast}将音频转为频谱图后用Vision Transformer处理，在多个音频分类基准上取得突破，但需大规模预训练数据。考虑到本工作的数据规模和可解释性需求，采用经典声学特征（MFCC + Mel频谱 + 统计量）+ 确定性投影的方案，兼顾效果与可解释性。

\section{数据增强与半合成方法}

在医疗低频事件（如跌倒）场景中，真实数据稀缺且涉及隐私。特征级数据增强是缓解数据不足的有效方法。Mixup\cite{zhang2018mixup}通过对样本对进行线性插值生成新样本，标签同步软化，能起到正则化作用。SMOTE等过采样技术也可在特征空间扩充少数类样本。

本工作采用``真实特征 + 场景规则组合 + 特征级增强''的半合成策略：从真实数据提取特征原型，按场景规则组合（跌倒=真实跌倒视频特征+真实撞击音频特征+异常生理特征），并对半合成样本加入小幅高斯噪声增强多样性。这种方法在保证特征真实性的同时扩充了数据规模。

\section{与本工作的差异定位}

现有老年人监护系统多采用固定权重融合或单一模态深度学习。本工作的差异在于：

\begin{enumerate}
    \item \textbf{可学习的跨模态注意力}：采用4层Cross-Modal Attention实现四模态联合预测，而非固定权重加权；
    \item \textbf{四模态真实数据}：四模态全部基于真实公开数据特征训练（ESC-50音频 + Pexels视频），而非模拟噪声；
    \item \textbf{系统性消融验证}：通过三组消融实验量化各模态真实贡献，证明四模态均参与决策；
    \item \textbf{面向部署的实用性}：集成缺失数据处理（默认特征填充）与可解释性展示（模态权重+注意力热力图），开发完整Web应用。
\end{enumerate}
